\section{Vorstellung der ausgewählten Fallstudie}

In diesem Abschnitt soll die ausgewählte Fallstudie vorgestellt werden, die im Rahmen dieser Arbeit analysiert wird.

\subsection{Auswahl der Studie}

Welche Studie dieser Arbeit eine geeignete und ausreichende Grundlage zur Analyse der Triangulation in der Evaluation von Präventionsmaßnahmen gegen Kinderarmut ist, hängt von nachfolgenden Kriterien ab. Es sollte sich um eine Fallstudie handeln, die einen Mixed-Methods-Ansatz verfolgt, einem typsischen Studiendesign für methodische Triangulation. Methodik und Erkenntnisgewinn beider Foschungsparadigmen sollten dabei in einem ausreichenden Detailgrad beschrieben worden sein, um sie in dieser Arbeit analysieren zu können. Darüber hinaus sollte sie sich inhaltlich auf die Themenwahl der Kinderarmutsprävention beziehen, um die Vorteile der Triangulation in eben diesem Kontext herausarbeiten zu können.

Ausgewählt wurde die Studie: Gute Chancen für alle Kinder (2023) im Rahmen der Evaluation des Tübinger Präventionskonzeptes gegen Kinderarmut. Die Studie trifft den thematischen Schwerpunkt der Kinderarmutsprävention mit konkreten und umfassend dokumentierten Maßnahmen der Stadt Tübingen. Das Foschungsdesign entspricht einem Mixed-Methods-Ansatz mit einer ausführlichen Beschreibung der verschiedenen Methoden. Schließlich wird im Sinne der Auftraggeber ein großer Schwerpunkt auf den Erkenntnisgewinn gelegt, der eine Analyse des selbigen im Sinne der Triangulation für diese Arbeit eröffnet.

\subsection{Kontext und Forschungsgegenstand der Studie}

2014 gründeten die Stadt Tübingen, die Liga der freien Wohlfahrtspflege und das Bündnis für Familie den Runde Tisch gegen Kinderarmut, der sich mit Angeboten und Maßnahmen auseinandersetzt, welche die Teilhabe und Entwicklungsmöglichkeiten der von Armut betroffenen Kinder und Familien in Tübingen stärken soll. Daraus formte sich das Angebots- und Maßnahmenprogramm des Tübinger Präventionskonzepts gegen Kinderarmut. In der ausgewählen Studie Gute Chancen für alle Kinder wurde das Tübinger Präventionskonzept gegen Kinderarmut von der Evangelischen Hochschule Ludwigsburg auf Anfrage und in Kooperation mit der Stadt Tübingen evaluiert.

\subsection{Zielsetzung und Fragestellung der Studie}

Die Studie hat das Ziel der Evaluation der Angebote und Maßnahmen des Tübinger Präventionskonzeptes und inwiefern die beabsichtigte Wirkungsweise erzielt werden konnte. Konkret unterscheiden die Autoren drei Forschungsfragen. Die erste Forschungsfrage fragt nach bewirkten Veränderungen, den Erfolgen und den bestehenden Hürden. Die zweite Forschungsfrage fragt, ob die durch die Maßnahmen beabsichtigte Wirkung erzielt werden konnte. Dabei soll insbesondere darauf eingegangen werden, wie bekannt die Angebote sind, wie gut über sie informiert wird und wie hoch ihre Annahme durch die Zielgruppe ist. Schließlich sollen auch weitere Bedarfe von Armut betroffener oder bedrohter Kinder und Familien erhoben werden.

\subsection{Forschungsdesign der Studie}

Das Studiendesign besteht aus einer quantitativen Erhebung und mehreren qualitativen Methoden. Die Kombination aus quantitativen und qualitativen Methoden macht es zu einem Mixed-Methods-Ansatz. Anders als einfache Mixed-Methods-Ansätze bei denen qualitative Methoden häufig nur vor quantitativen Methoden zur Exploration oder nach quantitativen Methoden zur Vertiefung eingesetzt werden, enthält die Studie sowohl Exploration zu Beginn als auch Vertiefung zum Schluss.

Wie in Abbildung 1 TODO zu sehen, wurde mit einer Literatur- und Datenanalyse begonnen, um einen Überblick über die Angebote des Präventionskonzepts zu gewinnen. Die Erkenntnisse daraus wurden für die Entwicklung eines Leitfadens für die Fokusgruppe aus Haupt- und Ehrenamtlichen genutzt. Die Erkenntnisse daraus wurden wiederum zur Fragebogenkonstruktion für die Online-Befragung der von Armut betroffenen oder bedrohten Familien. Die Ergebnisse wurden in Fokusgruppen aus Eltern aus dieser Zielgruppe vertieft und schließlich wurden die gesammelten Ergebnisse in einem Workshop des Runden Tisches reflektiert.

Nach der Design-Notation von Morse (Morse, 1991 TODO) ließe sich das Mixed-Methods-Design hier als "qual -> qual -> QUAN -> QUAL -> qual" beschreiben. Hervorzuheben ist dabei, dass die Methoden nicht parallel sondern stets aufeinander aufbauend eingesetzt wurden. Kernstücke bilden die quantitative Online-Befragung und anschließenden Fokusgruppen der von Armut betroffenen oder bedrohten Familien.

% TODO explanatives Design (https://www.socialnet.de/lexikon/Empirische-Sozialforschung#toc_4) und konvergente/divergente/komplementäre Ergebnisse